\section{Unequal precision in branch observations}
\label{sec:unequal-branches}

This section proves a second, unconditional extension of the source article.
Question \texttt{lub:q:adversarial} of \cite{article} asks for the branch-recovery results to be extended
to unequal truncation lengths and to coefficients modulo higher powers of a
uniformizer.  We resolve the unequal-length part over a field, including the
exact observation kernel, the complete error/erasure criterion, an exact
contact-tree algorithm, and optimal total recording budgets.  Coefficients in
$\mathcal O/\pi^r$ for $r>1$ are outside the theorem.

\subsection{The kernel and the recoverable jet}

Let $F$ be a field, put $R=F[[X]]$ and $\mathfrak m=(X,Y)$, and let
$u_b\in XR$, $b\in I$, be $N$ pairwise distinct branches.  Set
\[
 s_{bc}=\operatorname{ord}_X(u_b-u_c)\ge1\quad(b\ne c).
\]
Choose arbitrary integers $n_b\ge0$.  A branch with $n_b=0$ has an empty record.
The observations of an otherwise unrestricted $H\in F[[X,Y]]$ are
\[
 \Omega_{\boldsymbol n}(H)_b
   =H(X,u_b(X))\bmod X^{n_b}.
\]
The target is the complete total-degree jet
$J_d(H)=H\bmod\mathfrak m^d$, not merely the coefficients of a polynomial in $Y$.
It has $d(d+1)/2$ scalar coefficients.

Order the branches as follows.  Having chosen $b_0,\ldots,b_{j-1}$, choose an
unused $b_j$ maximizing
\begin{equation}
 \mu_j=n_{b_j}-\sum_{i<j}s_{b_jb_i},
 \qquad
 \Pi_j(Y)=\prod_{i<j}(Y-u_{b_i}(X)).
 \label{eq:unequal-greedy}
\end{equation}
All ties are allowed.  The scores may be negative; write
$\mu_j^+=\max(0,\mu_j)$.  Every unused score falls by at least one at the next
step, so
\begin{equation}
 \mu_{j+1}\le\mu_j-1.
 \label{eq:unequal-descent}
\end{equation}

\begin{theorem}[Exact kernel at unequal precisions]
\label{thm:unequal-kernel}
As an $R$-module, the observation kernel is
\begin{equation}
 \ker\Omega_{\boldsymbol n}
 =\Pi_N F[[X,Y]]\ \oplus\!
   \bigoplus_{j=0}^{N-1}X^{\mu_j^+}R\Pi_j.
 \label{eq:unequal-kernel}
\end{equation}
For $1\le d\le N$, the records determine $J_d(H)$ for every $H$ if and only if
\begin{equation}
 \mu_{d-1}>0.
 \label{eq:unequal-clean}
\end{equation}
The largest completely recoverable jet order is the number of positive scores.
Moreover, the dimension of the unresolved part of the target is exactly
\begin{equation}
 \dim_F\frac{\ker\Omega_{\boldsymbol n}+\mathfrak m^d}{\mathfrak m^d}
 =\sum_{j=0}^{d-1}\max\{0,d-j-\mu_j^+\}.
 \label{eq:unequal-dimension}
\end{equation}
The nonnegative exponents $\mu_j^+$ are independent of the choices at ties.
\end{theorem}

\begin{proof}
Successive division by $Y-u_{b_i}(X)$ gives a unique expansion
\[
 H=\Pi_N Q+\sum_{j<N}c_j(X)\Pi_j(Y),\qquad c_j\in R.
\]
This division is valid for formal series: $u_b\in XR$ makes substitution
well defined, and
$Y^a-u_b^a=(Y-u_b)\sum_{t=0}^{a-1}Y^{a-1-t}u_b^t$ gives the quotient term by
term in the $\mathfrak m$-adic topology.  Uniqueness of the remainder follows
by evaluating successively at $u_{b_0},u_{b_1},\ldots$; the relevant products
of distinct branch differences are nonzero in the domain $R$.

For each $j$, the function $X^{\mu_j^+}\Pi_j$ satisfies every observation
constraint.  It vanishes identically on the previously selected branches.
At any remaining branch $b$, the maximality in
\eqref{eq:unequal-greedy} gives
\[
 \mu_j^+ +\sum_{i<j}s_{bb_i}\ge n_b.
\]
Conversely suppose the remainder is in the kernel.  Evaluation at $b_0$
forces $v_X(c_0)\ge n_{b_0}=\mu_0$.  Subtract this term, which itself lies
in the kernel.  At step $j$ the previously removed terms satisfy all
constraints and the later Newton factors vanish at $b_j$, so
\[
 v_X(c_j)+\sum_{i<j}s_{b_jb_i}\ge n_{b_j}.
\]
Together with $c_j\in R$ this is $v_X(c_j)\ge\mu_j^+$.  The two implications
prove \eqref{eq:unequal-kernel}.

The Newton change of basis respects total degree.  The coefficient of $Y^i$
in $\Pi_j$, for $i<j$, belongs to $X^{j-i}R$, whereas $\Pi_j$ is monic of
$Y$-degree $j$.  Consequently
\[
 \operatorname{ord}_{\mathfrak m}\Bigl(\sum_{j<N}c_j\Pi_j\Bigr)
   =\min_{j<N}\{v_X(c_j)+j\}.
\]
For the converse inequality in this equality, among the summands of least
total order take the one with largest $Y$-degree: its leading coefficient
cannot cancel with any of the other leading forms.  Also
$\Pi_N\in\mathfrak m^N$.

If $\mu_{d-1}\ge1$, \eqref{eq:unequal-descent} gives
$\mu_j+j\ge\mu_{d-1}+d-1\ge d$ for every $j<d$; all remaining factors already
have total order at least $d$.  The kernel therefore lies in
$\mathfrak m^d$.  If $\mu_{d-1}\le0$, the kernel contains $\Pi_{d-1}$, whose
total order is exactly $d-1$, proving failure.  The classes of
$X^i\Pi_j$ with $i+j<d$ form a basis of the target jet, so counting the
kernel basis classes gives \eqref{eq:unequal-dimension}.

Finally, within the $R$-module of kernel polynomials of $Y$-degree at most
$j<N$, the possible leading coefficients form the intrinsic ideal
$X^{\mu_j^+}R$.  This ideal does not depend on the ordering and proves
tie independence of the capped scores.
\end{proof}

When all $n_b=n$, the rule is ordinary $P$-ordering, and
$\mu_j=n-\lambda_j$, where $\lambda_j$ is the greedy contact sequence of the
source article.  Thus the theorem specializes exactly to
\texttt{lub:thm:kernel}.  The weights $n_b$ cannot be replaced by their
minimum or average without losing information.

\subsection{A sparse certificate and the exact recording budget}

\begin{corollary}[Only $d$ clean branches are needed]
\label{cor:unequal-certificate}
The full observations recover $J_d$ if and only if some $d$-element subset
$S\subseteq I$ satisfies
\begin{equation}
 n_b\ge1+\sum_{c\in S\setminus\{b\}}s_{bc}\qquad(b\in S).
 \label{eq:unequal-certificate}
\end{equation}
When recovery holds, the first $d$ branches in
\eqref{eq:unequal-greedy} give such a subset.
\end{corollary}

\begin{proof}
The first $d$ choices remain a greedy order when the process is restricted
to that $d$-element set, with the same first $d$ scores.  By
Theorem~\ref{thm:unequal-kernel}, its observations recover $J_d$ whenever
those of the whole set do.

For exactly $d$ branches put
$\Delta_b=\prod_{c\ne b}(u_b-u_c)$ and $\sigma_b=v_X(\Delta_b)$.
Lagrange interpolation of the remainder of $H$ modulo
$\prod_{b\in S}(Y-u_b)$ gives
\[
 R_H(Y)=\sum_{b\in S}\frac{H(X,u_b)}{\Delta_b}
                  \prod_{c\in S\setminus\{b\}}(Y-u_c).
\]
If all records vanish and $n_b\ge\sigma_b+1$, each quotient is in $XR$;
hence every summand belongs to $\mathfrak m^d$.  If $n_b\le\sigma_b$ for
some $b$, the explicit function
\[
 \frac{X^{\sigma_b}}{\Delta_b}
       \prod_{c\in S\setminus\{b\}}(Y-u_c)
\]
is integral, has a nonzero jet of order $d$, vanishes on the other branches,
and has value $X^{\sigma_b}$ on $b$.  It is indistinguishable from zero.
This proves the necessary and sufficient componentwise criterion.
\end{proof}

One may choose the precisions before observing $H$.  The recording cost is
$\sum_b n_b$; unrecorded branches have precision zero.

\begin{theorem}[Minimum cost of clean jet recovery]
\label{thm:unequal-budget}
For fixed distinct branches and $1\le d\le N$, the least total recording
cost that recovers $J_d$ is
\begin{equation}
 B_d=d+2\min_{\substack{S\subseteq I\\|S|=d}}
                    \sum_{\{b,c\}\subseteq S}s_{bc}.
 \label{eq:unequal-budget}
\end{equation}
If $\lambda_j$ is the ordinary minimum-contact $P$-ordering sequence on all
branches, this is also
\[
 B_d=d+2\sum_{j=0}^{d-1}\lambda_j.
\]
Choose the first $d$ branches of an ordinary $P$-ordering, give each the
precision in the right side of \eqref{eq:unequal-certificate}, and give all
other branches precision zero.  This attains the minimum.
\end{theorem}

\begin{proof}
Corollary~\ref{cor:unequal-certificate} supplies a $d$-element set $S$ for
any successful profile, and summing its necessary componentwise bounds gives
$\sum_b n_b\ge d+2\sum_{\{b,c\}\subseteq S}s_{bc}$.
The same componentwise bounds show that the profile constructed from any
minimizing $S$ succeeds with exactly that cost.

For completeness, ordinary $P$-ordering attains the displayed pair-energy
minimum.  Let $a_0,a_1,\ldots$ be such an ordering, let
$P_j(Y)=\prod_{i<j}(Y-u_{a_i})$, and let $\alpha_j=P_j(u_{a_j})$, so
$v_X(\alpha_j)=\lambda_j$.  By its greedy definition every
$P_j/\alpha_j$ is integral valued on every branch.  For any $d$-element set
$S$, the evaluation matrix of $1,P_1/\alpha_1,\ldots,P_{d-1}/\alpha_{d-1}$
on $S$ has integral entries.  Its determinant is a Vandermonde determinant
divided by $\prod_{j<d}\alpha_j$, whence
\[
 \sum_{\{b,c\}\subseteq S}s_{bc}\ge\sum_{j<d}\lambda_j.
\]
For $S=\{a_0,\ldots,a_{d-1}\}$ the matrix is triangular with diagonal one,
so equality holds.
\end{proof}

\subsection{Arbitrary corrupted or erased branch records}

The block distance between distinct target jets is
\[
 \delta_d(\boldsymbol n;I)
   =\min_{H\notin\mathfrak m^d}
          \#\{b:\Omega_{\boldsymbol n}(H)_b\ne0\}.
\]
For a subset $S$, rerun \eqref{eq:unequal-greedy} with its inherited
precisions and write its scores as $\mu_j(S)$, setting
$\mu_j(S)=-\infty$ for $j\ge|S|$.  Define
\[
 C_I(j,T)=\max\{|S|:S\subseteq I,\ \mu_j(S)\le T\}.
\]
For the distance formula only $T=0$ is needed.

\begin{theorem}[Exact error and erasure threshold]
\label{thm:unequal-errors}
For $1\le d\le N$,
\begin{equation}
 \delta_d(\boldsymbol n;I)=N-C_I(d-1,0).
 \label{eq:unequal-distance}
\end{equation}
The target jet is uniquely determined for every $H$ despite at most $w$
arbitrarily corrupted blocks and at most $z$ flagged erased blocks if and
only if
\begin{equation}
 2w+z<\delta_d(\boldsymbol n;I).
 \label{eq:unequal-errors}
\end{equation}
All failures have explicit witnesses in $F[[X]][Y]$ of $Y$-degree at most
$d-1$ and total order at most $d-1$.  Finite polynomial witnesses result
by truncating their $X$-coefficients sufficiently far.
\end{theorem}

\begin{proof}
If $|S|\ge d$, Theorem~\ref{thm:unequal-kernel} gives a nonzero target jet
whose observations vanish on $S$ exactly when $\mu_{d-1}(S)\le0$.
In this case $\Pi_{d-1}$ for the greedy order on $S$ is a witness.
For $|S|<d$, the product $\prod_{b\in S}(Y-u_b)$ is a witness at every
precision.  Thus $C_I(d-1,0)$ is exactly the largest possible zero set of
an observation word with nonzero target, proving
\eqref{eq:unequal-distance}.

Two explanations of the same received records differ on at most the
union of their error locations and the erased locations, namely $2w+z$
blocks.  If this is smaller than the target distance, their target jets
coincide.  Conversely choose a minimum-distance word with nonzero target
and compare it with zero.  If $\delta_d\le2w+z$, erase up to $z$ of its
nonzero blocks and partition the remaining at most $2w$ blocks into two
sets of size at most $w$.  On one set let the received word agree with
zero, and on the other let it agree with the chosen word.  Both target
jets then explain the data.  When $\delta_d=0$, the ambiguity exists even
without errors.  To obtain a finite polynomial witness, truncate each
$X$-coefficient at $M=\max\{d,n_{\max}\}$, where
$n_{\max}=\max_b n_b$.  The discarded terms belong to $X^MF[[X,Y]]$;
they affect neither the target jet nor any record.
\end{proof}

This is a target-jet distance: words caused solely by terms of total degree
at least $d$ need not have positive distance.  It is therefore not legitimate
to replace it by the minimum distance of the entire observation code.

\begin{corollary}[Exact componentwise saturation and its cost]
\label{cor:unequal-saturation}
Always $\delta_d\le N-d+1$.  Equality holds if and only if, for each $b$,
\begin{equation}
 n_b\ge 1+
     \text{the sum of the $d-1$ largest numbers $s_{bc}$, $c\ne b$}.
 \label{eq:unequal-saturation}
\end{equation}
The least total recording cost attaining distance $N-d+1$ is the sum of
these $N$ componentwise lower bounds.  In particular, if
$N=d+2w+z$, this is exactly the minimum cost of correcting $w$ corrupted
blocks and $z$ erased blocks.
\end{corollary}

\begin{proof}
A product of $d-1$ branch equations has a nonzero target and vanishes on
those branches, proving the upper bound.  Equality is equivalent to
clean recovery from every $d$-element subset.  Apply
\eqref{eq:unequal-certificate} to each such subset and maximize the sum
of contacts at each fixed $b$.  The resulting inequalities are independent
componentwise lower bounds and attain equality simultaneously.  When
$N=d+2w+z$, condition \eqref{eq:unequal-errors} requires the maximum
possible distance $N-d+1$.
\end{proof}

\subsection{The exact contact-tree algorithm}

The contact numbers satisfy
$s_{bc}\ge\min\{s_{ba},s_{ac}\}$.  Represent them by a rooted tree whose
leaves are the branches: the common path length of two distinct leaves is
their contact.  The root has a stem of length $\min s_{bc}$; a nonleaf
cluster splits according to strict increase of contact.  Let $\ell_v$ be
the incoming length of a node, let $I_v$ be its leaves, and use the residual
contact lengths below its parent.  At a leaf put $\mu_0^v=n_b$ and
$\mu_j^v=-\infty$ for $j\ge1$.

\begin{lemma}[Weighted merge]
\label{lem:unequal-merge}
Let $z_0\ge z_1\ge\cdots$ be the decreasing rearrangement of the union of
the finite score lists of the children of $v$.  Then
\[
 \mu_j^v=z_j-j\ell_v\qquad(0\le j<|I_v|).
\]
In particular the full score list, including its negative entries, depends
only on the contact tree and the precision labels, and not on tie choices.
\end{lemma}

\begin{proof}
After $j$ choices every remaining candidate pays the same $j\ell_v$
for the common stem.  Its remaining penalty involves only choices in its
own child.  Each child offers a nonincreasing stream of scores, by
induction from the leaves.  The maximum rule therefore merges those
streams in decreasing order and subtracts $j\ell_v$ at step $j$.
\end{proof}

Define
\[
 C_v(j,T)=\max\{|S|:S\subseteq I_v,\ \mu_j^v(S)\le T\},
\]
again using $-\infty$ beyond the finite score list.

\begin{theorem}[Pruning recursion]
\label{thm:unequal-tree}
If $j\ge|I_v|$ or $T\ge\max_{b\in I_v}n_b$, then $C_v(j,T)=|I_v|$.
For a leaf $b$,
\[
 C_b(0,T)=\begin{cases}1,&n_b\le T,\\0,&n_b>T,\end{cases}
 \qquad C_b(j,T)=1\quad(j\ge1).
\]
At an internal node,
\begin{equation}
 C_v(j,T)=
 \max_{\substack{j_c\ge0\\\sum_c j_c\le j}}
       \sum_c C_c(j_c,T+j\ell_v).
 \label{eq:unequal-pruning}
\end{equation}
Backtracking maximizing allocations gives a maximal zero set and therefore
a sharp ambiguity witness in Theorem~\ref{thm:unequal-errors}.
For scores up to $d-1$, a direct integer-threshold implementation uses
$O(Nd^2(n_{\max}+1))$ elementary arithmetic operations, where
$n_{\max}=\max_b n_b$; all larger thresholds can be capped.
\end{theorem}

\begin{proof}
By Lemma~\ref{lem:unequal-merge}, including the $-\infty$ tail,
$\mu_j^v(S)\le T$ is equivalent to the $j$th entry of the merged child
streams being at most $T+j\ell_v$.  This says that at most $j$ entries
of those streams are strictly greater than that threshold.  If child
$c$ contributes $h_c$ such entries, then $\sum_c h_c\le j$, and
$h_c\le j_c$ is exactly the condition
$\mu_{j_c}^c(S\cap I_c)\le T+j\ell_v$.  This gives both directions of
\eqref{eq:unequal-pruning} by taking maximizing sets in the children.
The base cases are immediate.  Starting at $T=0$, thresholds never
become negative; a threshold at least $n_{\max}$ imposes no restriction.
There are $O(N)$ nodes in the compact contact tree, $d$ ranks and at
most $n_{\max}+1$ relevant integer thresholds.  For each child threshold $U$, combine the child tables for all rank
budgets simultaneously by max-convolution, at cost $O(d^2)$ per child.
The table entry $(j,T)$ at the parent then reads the resulting rank-$j$
entry at $U=\min\{n_{\max},T+j\ell_v\}$.  This proves the operation bound.
\end{proof}

The algorithm computes the exact distance and a worst surviving
configuration.  It is not a polynomial-time decoder for corrupted records:
one may still decode by enumerating candidate error sets and solving finite
linear systems, and the faster error-location part of the source question
remains open here.

\subsection{Two exact examples and verification}

Over $F=\mathbb F_2$, take $u_b=(1+X)^b-1$ for $0\le b\le7$ and precisions
\[
 (n_0,\ldots,n_7)=(6,4,9,5,8,3,7,4).
\]
The contacts are $s_{bc}=2^{v_2(b-c)}$.  Breaking ties by the smallest
index, the greedy order and scores are
\[
 (2,4,3,1,6,0,7,5),\qquad(9,6,3,0,-1,-4,-6,-9).
\]
Thus the $46$ recorded coefficients recover the complete order-three jet,
consisting of the six coefficients of $1,X,Y,X^2,XY,Y^2$, but do not
recover the order-four jet.  The exact order-three distance is $4$, so
one arbitrary corrupt record together with one flagged erased record
is correctable.  Two arbitrary corruptions cannot always be corrected.
A sharp nonzero difference is
\[
 H(X,Y)=Y(Y+X+X^2+X^3).
\]
Its records vanish on branches $0,3,5,7$ and are nonzero precisely on
$1,2,4,6$.

For an exact allocation example, work over $\mathbb F_5$ with branches
\[
 0,\quad X^5,\quad X,\quad2X,\quad3X.
\]
The first two have contact $5$; every other pair has contact $1$.
For the order-three jet, the least componentwise profile giving maximal
distance $3$ is
\[
 (7,7,3,3,3).
\]
Its total cost is $23$, and it corrects one arbitrary corrupted branch.
Every successful profile for this five-branch, one-corruption problem has
cost at least $23$.  A common precision needs $7$ coefficients at every
branch, costing $35$.

The accompanying script \texttt{verify\_unequal\_branches.py} verifies the
claims using exact finite-field arithmetic.  It exhausts all $5^4=625$
precision profiles with values $0,\ldots,4$ on four binary multiplicative
branches; adds $120$ deterministic random configurations over
$\mathbb F_2,\mathbb F_3,\mathbb F_5$; and checks the two displayed examples.
Across these $747$ profiles it performs $2{,}996$ comparisons of the
unresolved-jet formula against direct expansion and matrix rank,
$2{,}996$ comparisons of the tree recursion against exhaustive branch
subsets, and $2{,}996$ checks of the componentwise saturation criterion.
These finite checks corroborate the proof; they are not substitutes for it.

\subsection{Attribution and scope of novelty}

The use of Newton polynomials and $P$-orderings is classical
\cite{bhargava}.  Weighted greedy selection on an ultrametric is likewise
established: Grinberg and Petrov \cite{grinbergpetrov} treat element weights
and ultrametric pair contributions.  The rule here is their greedy
perimeter rule with weight $n_b$ and pair distance $-s_{bc}$.
Amini \cite{amini} supplies the related unweighted tree-factorial framework.

The contribution claimed here is the exact unequal-truncation observation
kernel for unrestricted bivariate formal series, its total-degree target
interpretation, and the resulting sharp deletion/error and recording-budget
formulas.  These resolve the specified unequal-length question in the
supplied article.  A targeted literature search found no prior statement of
this combined recovery result, but that is not a proof of worldwide priority.
The seven ZIP reports in the repository's incoming directory at the pinned
commit were also checked; none states this extension.  No assertion is made
that $P$-orderings, weighted greedy optimality, Lagrange interpolation, or the
standard distance inequality $2w+z<\delta$ are new.
